\chapter*{\rm\bfseries Abr\'eg\'e}
\addcontentsline{toc}{chapter}{Abr\'eg\'e}

La saisie robotis\'ee de tissus minces et d\'eformables demeure un
probl\`eme ouvert en automatisation industrielle. Deux g\'eom\'etries de
pr\'ehenseur sont couramment propos\'ees~: les doigts rigides \`a
m\^achoires parall\`eles, qui offrent un contr\^ole de force pr\'ecis
mais une tol\'erance op\'erationnelle \'etroite, et les doigts souples
de type finray, qui r\'eduisent la pr\'ecision de force pour une
enveloppe op\'erationnelle plus large. [Cette th\`ese pr\'esente une
\'evaluation exp\'erimentale de ce compromis pour saisir un
tissu plat pos\'e sur une surface, en comparant deux configurations
rigides \`a m\^achoires parall\`eles et quatre configurations finray
imprim\'ees en 3D, soit six variantes de pr\'ehenseur, sur un bras
Franka Research 3.] La profondeur d'appui et l'angle d'inclinaison ont \'et\'e
choisis comme axes d'\'evaluation parce qu'ils repr\'esentent les
erreurs de placement qu'un syst\`eme d\'eploy\'e ne peut pas totalement
\'evaluer: un bras peu pr\'ecis ou grossi\`erement calibr\'e porte
une erreur de hauteur de surface \`a l'\'echelle millim\'etrique et une
erreur d'orientation \`a l'\'echelle du degr\'e, que le pr\'ehenseur
doit mitiger. La force de contact stationnaire a servi de variable de
r\'esultat principale et la r\'eussite de la prise de r\'esultat
secondaire.

Sur 840 essais analys\'es, les m\^achoires rigides n'ont tenu le tissu
que dans une fen\^etre de profondeur de $1{,}5$\,mm, ou de $2{,}5$\,mm
avec un coussinet souple, born\'ee par une limite de s\'ecurit\'e en
force. Sous inclinaison, ces fen\^etres se sont encore r\'etr\'ecies,
entre $0{,}5$ et $2{,}5$\,mm, et deux des conditions de la m\^achoire
d'origine n'ont produit aucune prise r\'eussie \`a aucune profondeur.
Cependant, chaque variante finray a conserv\'e toute la fen\^etre
d'analyse calibr\'ee de $10$\,mm \`a chaque orientation test\'ee, de
sorte que l'inclinaison n'a pas r\'eduit la plage utile du finray. En
outre, les quatre variantes finray suivent une seule courbe
force-profondeur une fois normalis\'ee par leur propre force maximale,
avec une erreur moyenne par paire inf\'erieure \`a $1\%$. La g\'eom\'etrie
du doigt fixe donc l'\'echelle de la force de contact, bien qu'elle ne
change pas la forme de la r\'eponse. Il est \`a noter que la force de
contact pr\`es de la limite de la plage utile varie encore avec le sens
de l'inclinaison, bien que cette sensibilit\'e dispara\^isse une fois le
doigt appuy\'e au-del\`a du coude de flambement.

\TODO{Faire v\'erifier cette traduction par un coll\`egue francophone
avant le d\'ep\^ot final (exigence de qualit\'e, non de conformit\'e~:
GPS exige seulement la pr\'esence d'un abr\'eg\'e en fran\c{c}ais).}
